\documentclass{article}
\usepackage{graphicx} % Required for inserting images

\usepackage{amsmath}
\usepackage{amssymb}

\title{Combinatorics 3}
\author{Siddharth K}
\date{October 2025}

\begin{document}

\maketitle

\section*{Problem 1}
With PIE:

Let $A_i$ be the set of such functions where $f(i) = i$. Then we want to count
$$\left|\bigcap_{i \in [n]} \overline{A_i}\right| = n^n - \left|\bigcup_{i\in [n]} A_i\right|.$$
We have that $|A_i| = n^{n - 1}$, since we can fix one element, $i$, and the remaining can map anywhere. Extending this, we have that $|A_{i_1} \cap \cdots \cap A_{i_k}| = n^{n - k}$. The number of ways to choose the $k$ indices is $\binom{n}{k}$, so using PIE,
$$n^n - \left|\bigcup_{i\in [n]} A_i\right| = n^n - \sum_{k = 1}^n (-1)^{k - 1}\binom nk n^{n - k} = \sum_{k = 0}^n (-1)^{k}\binom nk n^{n - k}.$$

Without PIE:

We assign each element of $[n]$ to any other element of $[n]$ except itself, giving $(n - 1)^n$.

\section*{Problem 2}
We can count the number of ways to choose a hand of 13 cards that only contains 2-10 and subtract the result from $\binom{52}{13}$. Since there are $9 \cdot 4 = 36$ such cards, the result is $\binom{52}{13} - \binom{36}{13}$.

\section*{Problem 5}
If the excedance set contains $n - 1$, then $a_{n - 1} = n$. The number of such permutations is $(n - 1)!$. If the excedance set contains $n - 2$, then $a_{n - 2} = n - 1$ or $n$, so there are $2(n - 1)!$ such permutations. Finally, if it contains $n - 1$ and $n - 2$, then $a_{n - 2} = n - 1$ and $a_{n - 1} = n$, so there are $(n - 2)!$ permutations. By PIE, the answer is $3(n - 1)! - (n - 2)!$.

\section*{Problem 7}
Let $I(n)$ be the number of involutions on $[n]$. Every involution on $[n]$ either fixes $n$ or pairs it with some $i \in [n - 1]$. If it fixes $n$, the restriction to $[n - 1]$ is still an involution, so there are $I(n - 1)$ such involutions. If $n$ is paired with some $i$, there are $n - 1$ choices for $i$ and, after choosing $\{i, n\}$, the remaining $n - 2$ elements can be arranged in any involution, giving us $(n - 1)I(n - 2)$ possibilities. So
$$I(n) = I(n - 1) + (n - 1)I(n - 2),$$
with $I(0) = 1$ and $I(1) = 1$. We show that $I(n)$ is even by induction. We have that $I(2) = 2$ and $I(3) = 4$, which are even. Now, assume $I(m)$ is even for all $2 \le m \le n - 1$. For $n \ge 4$, $n - 1$ and $n - 2$ are $\ge 2$, so $I(n - 1)$ and $I(n - 2)$ are even. The recurrence gives $I(n) = I(n - 1) + (n - 1)I(n - 2)$, which is a sum of even numbers and is therefore even. So $I(n)$ is even, and the number of non-trivial involutions is $I(n) - 1$ which is odd.

\section*{Problem 8}
Let $X$ be finite and $f: X\to X$ a bijection. Then $f$ decomposes into disjoint cycles, so each $x \in X$ lies in a unique cycle $x \to f(x) \to \cdots \to f^{\ell_x}(x) = x$ for some positive integer $\ell_x$. Since $X$ is finite, there are finitely many cycles and cycle lengths, say $\ell_1, \ldots, \ell_k$. Let $N = \text{lcm}(\ell_1, \ldots, \ell_k)$. For any $x$, we have that $n$ is a multiple of $\ell_x$, so $f^n(x) = x$.

\section*{Problem 11}
Let $\pi$ be a derangement of $[n]$ and let $k = \pi(n)$. If $\pi(k) = n$, then $\pi$ corresponds to a derangement of $[n] \backslash \{n, k\}$, since $\{n, k\}$ is a cycle and the remaining $n - 2$ elements can be mapped freely. If $\pi(k) \neq n$, then create a derangement $\tau: [n - 1] \to [n - 1]$ by removing $n$ and replacing it with $k$. Specifically, $\tau(i) = \pi(i)$ if $\pi(i) \neq n$, and $\tau(i) = k$ if $\pi(i) = n$. The total number of ways to do this is $(n - 1)(d(n - 1) + d(n - 2))$. \\

We have that $d(1) = 0$, and using the formula, $d(1) = 1!(1 - 1) = 0$. Now, assume that $d(n) = n!\sum_{i = 0}^n \frac{(-1)^n}{i!}$. Then
$$d(n + 1) = n(d(n) + d(n - 1))$$
$$= n\left(n!\sum_{i = 0}^n \frac{(-1)^i}{i!} + (n - 1)!\sum_{i = 0}^{n - 1}\frac{(-1)^i}{i!}\right) = n\left(n!\sum_{i = 0}^n \frac{(-1)^i}{i!}\right) + n!\sum_{i = 0}^{n - 1}\frac{(-1)^i}{i!}$$
$$= (n + 1)\left(n!\sum_{i = 0}^n \frac{(-1)^i}{i!}\right) + (-1)^{n - 1} = (n + 1)!\sum_{i = 0}^{n + 1} \frac{(-1)^i}{i!},$$
as desired.

\end{document}